\documentclass[11pt,a4paper]{article}
\usepackage{auditstyle}
\lhead{Fresh Glimm--Jaffe proof audit}
\rhead{16 August 2026}
\title{Corrected structured proof I\\[3pt]
GJS Euclidean estimates $\Longrightarrow$ a cutoff-uniform local moment and physical Hamiltonian gap}
\author{Fresh reconstruction of \texttt{part\_i/main.tex}}
\date{16 August 2026}

\begin{document}
\maketitle

\begin{resultbox}
\textbf{Status.} \Conditional.  The derivation below is valid from the explicitly listed
cutoff-construction facts and the exact estimates in Glimm--Jaffe--Spencer (GJS),
Theorems 1.1.7--1.1.8.  It incorporates every local repair found in the audit: the
dimensionless-coupling reduction, the correct FKN references, signed truncations, the
two time-sign cases, the corrected unit-cell count, the missing exponential-moment
prefactor, and the restriction $0\leq\beta<m_1$.

This document proves no part of the GJS cluster expansion itself.  That theorem is an
import, displayed verbatim by page image in the separate citation dossier.
\end{resultbox}

\tableofcontents

\section{Objects and exact imports}

Let
\[
 \Hh=L^2(\mathcal S'(\R),d\phi_0),\qquad
 K(g)=H_0+V(g),\qquad H(g)=K(g)-E(g),
\]
where $E(g)=\inf\spec K(g)$ and
\[
 V(h)=\lambda\int_{\R}h(x){:}\mathscr P(\varphi_0(x)){:}\,dx,
 \qquad \mathscr P(\xi)=\sum_{n=0}^{2p'}a_n\xi^n,\quad a_{2p'}>0.
\]
The scalar product is conjugate-linear in its first argument.  The cutoffs satisfy
$0\leq g\leq1$ and have compact support.  Write $\Om_0=1$ for the free vacuum.

\begin{assumption}[cutoff Hamiltonian facts]\label{ass:cutoff}
For every such $g$, $K(g)$ is self-adjoint and bounded below, $E(g)$ is a simple
eigenvalue, and its normalized eigenvector $\Om_g$ can be chosen strictly positive
$d\phi_0$-almost everywhere.  Consequently
\[
 c_g:=\ip{\Om_g}{\Om_0}>0.
\]
For every compactly supported real $h$, the time-zero multiplication operator $V(h)$
is self-adjoint on its maximal domain.
\end{assumption}

\begin{assumption}[slab FKN input]\label{ass:fkn}
For a fixed $g$, the interacting Feynman--Kac--Nelson formula holds on every finite
time slab.  The correct source locations are Simon, Theorem III.6 for the free FKN
formula and Corollary V.14/Theorem V.15 for its interacting form; GRS, Theorems
II.14 and II.16, give the same transfer formula.  Simon's Theorem III.12, cited in
the audited manuscript, is not an FKN theorem.
\end{assumption}

Let $d\Phi$ be the Euclidean Gaussian measure with covariance
$(-\Delta+m_0^2)^{-1}$.  For compactly supported measurable $0\leq h\leq1$, put
\[
 dq_h=Z(h)^{-1}\exp\!\left[-\lambda\int h(z){:}\mathscr P(\Phi(z)){:}\,dz\right]d\Phi.
\]
GJS define localized sharp-time monomials $Q$ and a norm $\|Q\|_{\rm GJS}$ in
(1.1.7)--(1.1.9).  If $N(\Delta)$ is the total degree in a unit square $\Delta$, a
localized term with kernel $w$ has norm
\[
 K_1\left[\prod_\Delta K_2^{2N(\Delta)}N(\Delta)!\right]\|w\|_2,
\]
and a nonlocalized polynomial is normed by summing its localized pieces.

\begin{assumption}[CE1: GJS Theorem 1.1.8]\label{ass:ce1}
There is $C_8<\infty$, independent of the cutoff $h$, such that every admissible
sharp-time polynomial obeys
\[
 \left|\int Q(s)\,dq_h\right|\leq C_8\|Q(s)\|_{\rm GJS}.
\]
The theorem statement says ``uniformly as $h\uparrow1$''; GJS page 629 explicitly
states that the estimates of Theorem 1.1.8 are uniform in the space cutoff $h$.
\end{assumption}

\begin{assumption}[CE2: GJS Theorem 1.1.7]\label{ass:ce2}
If $Q_1,Q_2$ have disjoint localization squares separated by a strip of maximal
width $d$, then
\[
 \left|\int Q_1Q_2\,dq_h-\int Q_1\,dq_h\int Q_2\,dq_h\right|
 \leq C_7e^{-m_1d}\|Q_1\|_{\rm GJS}\|Q_2\|_{\rm GJS},                 \tag{CE2}
\]
where $m_1>0$ and $C_7$ are independent of $h,Q_1,Q_2$.
\end{assumption}

\begin{remark}[no circular gap assumption]
Assumptions \ref{ass:cutoff}--\ref{ass:ce2} include simplicity of the cutoff ground
state but do not include isolation of its eigenvalue.  The gap proved below is therefore
not used to obtain the slab limit.
\end{remark}

\section{The structured proof}

\begin{theorem}[Euclidean estimates imply (M) and (G)]\label{thm:EG}
There is a dimensionless threshold $\varepsilon_{\mathscr P}>0$ such that, if
$0\leq\lambda/m_0^2<\varepsilon_{\mathscr P}$, then the following statements hold
for every cutoff $g$.

\begin{enumerate}[label=\textup{(\roman*)}]
\item If $h$ is real, $\|h\|_\infty\leq1$, and $\supp h$ is contained in an
interval of length $2$, then
\[
 \|V(h)\Om_g\|\leq M,
 \quad
 M^2=9C_8K_1K_2^{8p'}(4p')!\lambda^2
       \left(\sum_{n=0}^{2p'}|a_n|\right)^2.                         \tag{M}
\]
\item
\[
 \spec H(g)\subseteq\{0\}\cup[m_1,\infty),                         \tag{G}
\]
and $0$ is a simple eigenvalue.
\end{enumerate}
\end{theorem}

\subsection*{Proof in Lamport form}

\LS{1}{1}{Reduce the GJS smallness hypothesis to $\lambda/m_0^2<\varepsilon_{\mathscr P}$.}

\LS{2}{1}{Fix source-admissible numbers $M_*>0$ and $\lambda_*>0$ such that the
GJS expansion is valid at mass $M_*$ whenever $0\leq\lambda'<\lambda_*$.}

\Because GJS formulate the construction for sufficiently large mass and sufficiently
small coupling.  Choosing one admissible pair merely names the constants in that
hypothesis.

\LS{2}{2}{For arbitrary $(m_0,\lambda)$ set $s=m_0/M_*$ and apply the unitary
space-time dilation stated on GJS page 590.  It sends}
\[
 (m_0,\lambda)\longmapsto (s^{-1}m_0,s^{-2}\lambda)
       =\left(M_*,\frac{M_*^2}{m_0^2}\lambda\right).                 \tag{2.1}
\]

\LS{2}{3}{The dilation introduces no change in the coefficients of $\mathscr P$ or
in Wick ordering.}

\Because In two dimensions the scalar field has scaling dimension zero.  More
explicitly, if $\widetilde\Phi(y)=\Phi(y/m_0)$, then
\[
 \E[\widetilde\Phi(y)\widetilde\Phi(y')]
 =C_{m_0}((y-y')/m_0)=C_1(y-y'),
\]
by the substitution $p=m_0q$ in the Fourier integral for $C_{m_0}$.  Equality of
all covariance contractions implies equality of the corresponding Wick polynomials.
The Jacobian contributes $m_0^{-2}$ to the interaction.

\LS{2}{4}{Set $\varepsilon_{\mathscr P}=\lambda_*/M_*^2$.  Then
$\lambda/m_0^2<\varepsilon_{\mathscr P}$ implies
$s^{-2}\lambda<\lambda_*$, so the transformed model satisfies the source
hypothesis.}

\LS{2}{5}{This proves $\langle1\rangle1$.}\QedStep

\LS{1}{2}{Establish the exact slab transfer identities without assuming a gap.}

Fix $g$ and define
\[
 h_{T,g}(t,x)=\mathbf1_{[-T,T]}(t)g(x),\qquad
 \psi_a=e^{-aH(g)}\Om_0,\qquad \widehat\psi_a=\psi_a/\|\psi_a\|.
\]

\LS{2}{1}{For bounded nonnegative time-zero functions $F_1,F_2$, FKN gives, for
$0\leq t\leq2T$ and $|s|\leq T$,}
\begin{align}
 \int F_1(\Phi(-t/2))F_2(\Phi(t/2))\,dq_{h_{T,g}}
 &=\frac{\ip{F_1\psi_{T-t/2}}{e^{-tH(g)}F_2\psi_{T-t/2}}}
         {\|\psi_T\|^2},                                           \tag{2.2}\\
 \int F_1(\Phi(s))\,dq_{h_{T,g}}
 &=\frac{\ip{\psi_{T-s}}{F_1\psi_{T+s}}}{\|\psi_T\|^2}.          \tag{2.3}
\end{align}

\Because The unshifted formula contains $K(g)$ and the same factor $e^{-2TE(g)}$
in numerator and partition function.  Multiplying every semigroup by the corresponding
$e^{aE(g)}$ replaces $K(g)$ by $H(g)$ and cancels that factor exactly.

\LS{2}{2}{Equations (2.2)--(2.3) hold for bounded signed functions.}

\Because Write $F=F^+-F^-$ and expand bilinearly.  Each of the four terms is a
finite nonnegative FKN integral, so linearity is legitimate.

\LS{2}{3}{They extend to the unbounded functions used below.}

\Because For a real multiplication function $F$, set
$F^{[M]}=(-M)\vee(F\wedge M)$.  Apply the bounded identity to $F^{[M]}$.
Whenever $F\psi_a\in\Hh$, $F^{[M]}\psi_a\to F\psi_a$ by dominated convergence,
because $|F-F^{[M]}|^2|\psi_a|^2\leq |F|^2|\psi_a|^2$.  On path space, use
$|F_1F_2|\leq(F_1^2+F_2^2)/2$ after translating both factors to their respective
time slices.  The CE1 estimates below supply the required square integrability.

\LS{2}{4}{This proves $\langle1\rangle2$.}\QedStep

\LS{1}{3}{Prove the slab vectors converge to the cutoff ground state using simplicity only.}

Let $P_g=|\Om_g\rangle\langle\Om_g|$.  Positivity gives
$P_g\Om_0=c_g\Om_g$ with $c_g>0$.  Hence
\[
 \psi_a=c_g\Om_g+e^{-aH(g)}(1-P_g)\Om_0.                            \tag{3.1}
\]

\LS{2}{1}{The second term of (3.1) converges to zero in norm.}

\Because If $\mu$ is the spectral measure of $(1-P_g)\Om_0$, then
\[
 \|e^{-aH(g)}(1-P_g)\Om_0\|^2
   =\int_{(0,\infty)}e^{-2aE}\,d\mu(E)\longrightarrow0
\]
by dominated convergence.  Simplicity says precisely that $\mu(\{0\})=0$; no
positive distance from zero is used.

\LS{2}{2}{Consequently}
\[
 \widehat\psi_a\to\Om_g,\qquad
 \|\psi_a\|\downarrow c_g,\qquad
 1\leq\frac{\|\psi_b\|}{\|\psi_a\|}\leq c_g^{-1}\quad(0\leq b\leq a), \tag{3.2}
\]
and every ratio in (3.2) tends to $1$ when both indices tend to infinity with
bounded difference.

\Because $e^{-aH(g)}$ is a contraction and $\|\Om_0\|=1$.

\LS{2}{3}{This proves $\langle1\rangle3$.}\QedStep

\LS{1}{4}{Derive the local interaction moment (M).}

Let $h$ satisfy the hypotheses of (M).  In Euclidean notation define
\[
 Q_V=\lambda^2\sum_{n,n'=0}^{2p'}a_na_{n'}
 \iint h(x)h(y){:}\Phi(0,x)^n{:}{:}\Phi(0,y)^{n'}{:}\,dx\,dy.       \tag{4.1}
\]

\LS{2}{1}{$Q_V$ is exactly the multiplication function $V(h)^2$.}

\Because Expanding the square of
$\lambda\sum_na_n\int h(x){:}\Phi(0,x)^n{:}dx$ gives (4.1).  No Wick
reordering of the product is performed or required; GJS (1.1.5) permits products
of individually Wick-ordered powers.

\LS{2}{2}{The support of $h$ meets at most three half-open unit spatial cells in
positive measure.}

\Because An interval of length $2$ meets at most three members of the partition
$[j,j+1)$ except possibly in endpoints, which have zero Lebesgue measure.

\LS{2}{3}{For one ordered pair of cells and degrees $n,n'$, the GJS local weight is
bounded by $K_1K_2^{8p'}(4p')!$.}

\Because The total local degree is at most $n+n'\leq4p'$.  If both variables
belong to one square, the weight is
$K_1K_2^{2(n+n')}(n+n')!$.  If they belong to two squares, it is
$K_1K_2^{2n}n!K_2^{2n'}n'!$.  In both cases
\[
 K_2^{2(n+n')}\leq K_2^{8p'},\qquad n!n'!\leq(n+n')!\leq(4p')!,
\]
after increasing $K_2$ to at least $1$, as GJS do.

\LS{2}{4}{Summing the at most $3^2=9$ ordered cell pairs and all coefficients gives}
\[
 \|Q_V\|_{\rm GJS}\leq
 9K_1K_2^{8p'}(4p')!\lambda^2
 \left(\sum_{n=0}^{2p'}|a_n|\right)^2=M^2/C_8.                     \tag{4.2}
\]

\LS{2}{5}{For every $T>0$,}
\[
 \|V(h)\widehat\psi_T\|^2
 =\int V(h)^2\,dq_{h_{T,g}}\leq C_8\|Q_V\|_{\rm GJS}\leq M^2.     \tag{4.3}
\]

\Because Apply (2.2) at $t=0$ with the two individually nonnegative functions
$F_1=F_2=|V(h)|$.  Their product is $V(h)^2=Q_V$.  Then apply CE1 and (4.2).
This is the required correction to the manuscript's invalid phrase
``$F_1F_2=V(h)^2\geq0$'', which did not specify nonnegative $F_1,F_2$.

\LS{2}{6}{Pass (4.3) to $\Om_g$.}

\Because Choose $T_k\to\infty$ such that
$\|V(h)\widehat\psi_{T_k}\|$ approaches its liminf.  By (4.3), a subsequence
converges weakly to some $y$.  The graph of the closed linear operator $V(h)$ is a
closed linear subspace of $\Hh\oplus\Hh$, hence is weakly closed.  Since
$\widehat\psi_{T_k}\to\Om_g$ strongly, $(\Om_g,y)$ lies in the graph.  Thus
$y=V(h)\Om_g$ and weak lower semicontinuity gives
\[
 \|V(h)\Om_g\|\leq\liminf_k\|V(h)\widehat\psi_{T_k}\|\leq M.
\]

\LS{2}{7}{This proves (M).}\QedStep

\LS{1}{5}{Obtain correct GJS norm bounds for polynomial time-zero vectors.}

Let
$Q=\prod_{i=1}^r\varphi_0(f_i)$ with real $f_i\in C_c^\infty(\R)$, all supported
in an interval $I$ of length $L$.  Let
\[
 N_I:=\#\{j:\ |I\cap[j,j+1)|>0\}.
\]

\LS{2}{1}{$N_I\leq\lfloor L\rfloor+2\leq L+2$.}

\Because If $j_{\min}$ and $j_{\max}$ are the extreme intersecting cells, then
$j_{\max}-j_{\min}-1<L$, hence the number
$j_{\max}-j_{\min}+1$ is at most $\lfloor L\rfloor+2$.  The manuscript's bound
$L+1$ is false for nonintegral $L$; an interval of length $0.2$ can meet two cells.

\LS{2}{2}{Decomposing each $f_i$ by the unit cells yields}
\begin{align}
 \|Q\|_{\rm GJS}
 &\leq K_1K_2^{2r}r!\,N_I^{r/2}\prod_{i=1}^r\|f_i\|_2,             \tag{5.1}\\
 \|Q^2\|_{\rm GJS}
 &\leq K_1K_2^{4r}(2r)!\,N_I^r\prod_{i=1}^r\|f_i\|_2^2
   =:\mathfrak n(Q)^2.                                             \tag{5.2}
\end{align}

\Because Expand the kernel into cell assignments.  For any assignment, the product
of local factorials is at most the factorial of total degree, and the $K_2$ exponents
sum to twice that degree.  For each $i$,
\[
 \sum_{j:I\cap[j,j+1)\ne\varnothing}\|f_i\mathbf1_{[j,j+1)}\|_2
 \leq N_I^{1/2}\|f_i\|_2
\]
by Cauchy--Schwarz.  Multiplying these $r$ estimates proves (5.1); applying it to
the $2r$ factors in $Q^2$ proves (5.2).

\LS{2}{3}{The norm is independent of the common sharp time $s$.}

\Because At every $s$ all factors occupy one temporal row.  Changing $s$ changes
only that common row label; the spatial cell assignments, local degrees, kernel norms,
and product of weights are unchanged.  This is the precise statement; arbitrary time
translation does not literally preserve the fixed lattice squares.

\LS{2}{4}{This proves $\langle1\rangle5$.}\QedStep

\LS{1}{6}{Transfer CE1 and CE2 to connected cutoff-ground-state semigroup bounds.}

\LS{2}{1}{CE1 and (2.2) at $t=0$ imply}
\[
 \sup_T\|Q\widehat\psi_T\|^2\leq C_8\mathfrak n(Q)^2.              \tag{6.1}
\]

\Because Apply the signed-truncation construction of $\langle2\rangle3$ to $Q$;
equivalently use $F_1=F_2=|Q|$.  Then the Euclidean integrand is $Q^2$, to
which CE1 and (5.2) apply.  The weak-graph argument of $\langle1\rangle4$ also gives
$\Om_g\in D(Q)$.

\LS{2}{2}{For fixed $s\in\R$,}
\[
 \int Q(s)\,dq_{h_{T,g}}\longrightarrow\ip{\Om_g}{Q\Om_g}.       \tag{6.2}
\]

\Because For $T>|s|$, (2.3) gives
\[
 \int Q(s)dq_{h_{T,g}}
 =\ip{\widehat\psi_{T-s}}{Q\psi_{T+s}/\|\psi_T\|}
   \frac{\|\psi_{T-s}\|}{\|\psi_T\|}.                            \tag{6.3}
\]
The first vector tends strongly to $\Om_g$, and the last factor tends to $1$.
Write the middle vector as
\[
 \frac{Q\psi_{T+s}}{\|\psi_T\|}
 =\frac{\|\psi_{T+s}\|}{\|\psi_T\|}Q\widehat\psi_{T+s}.          \tag{6.4}
\]
If $s\geq0$, the scalar ratio in (6.4) is at most $1$.  If $s<0$, it is at
most $c_g^{-1}$ by (3.2).  Thus (6.4) is bounded in either case.  Every weak
subsequential limit equals $Q\Om_g$ by weak closedness of the graph and because the
scalar ratio tends to $1$.  Hence the whole middle vector converges weakly to
$Q\Om_g$, proving (6.2).  This separate treatment of $s<0$ repairs the false
inequality $\|\psi_{T+s}\|\leq\|\psi_T\|$ used for both signs in the manuscript.

\LS{2}{3}{For $t\geq0$,}
\[
 \ip{Q\Om_g}{e^{-tH(g)}Q\Om_g}
 \leq\liminf_{T\to\infty}
 \int Q(-t/2)Q(t/2)\,dq_{h_{T,g}}.                                 \tag{6.5}
\]

\Because With
$u_T=Q\psi_{T-t/2}/\|\psi_T\|$, the signed-truncation extension of (2.2) gives
the exact equality
\[
 \int Q(-t/2)Q(t/2)dq_{h_{T,g}}=\|e^{-tH(g)/2}u_T\|^2.             \tag{6.6}
\]
The ratio bound (3.2) and (6.1) make $u_T$ bounded.  Write
$u_T=(\|\psi_{T-t/2}\|/\|\psi_T\|)Q\widehat\psi_{T-t/2}$; the scalar tends to
$1$ and is uniformly bounded.  From every subsequence choose a further one on
which $u_T$ converges weakly.  Since
$\widehat\psi_{T-t/2}\to\Om_g$ strongly, weak closedness of the graph of the
closed operator $Q$ identifies that weak limit as $Q\Om_g$.  Thus the full net
converges weakly to $Q\Om_g$.  Boundedness of $e^{-tH/2}$ and weak lower
semicontinuity of the norm prove (6.5).

\LS{2}{4}{For $t>2$, CE2 implies}
\[
 \left|\int Q(-t/2)Q(t/2)dq_h
 -\int Q(-t/2)dq_h\int Q(t/2)dq_h\right|
 \leq C_7e^{2m_1}e^{-m_1t}\|Q\|_{\rm GJS}^2.                      \tag{6.7}
\]

\Because Write the source-prescribed finite localization decomposition as
$Q=\sum_{a\in A}Q_a$, so that
$\|Q\|_{\rm GJS}=\sum_a\|Q_a\|_{\rm GJS}$.  The two temporal rows of every
pair $Q_a(-t/2),Q_b(t/2)$ are separated by a horizontal strip of width at least
$t-2$.  Therefore CE2 and the triangle inequality give
\[
 \begin{split}
 \left|\operatorname{Cov}_{dq_h}(Q(-t/2),Q(t/2))\right|
 &\le\sum_{a,b}\left|\operatorname{Cov}_{dq_h}
       (Q_a(-t/2),Q_b(t/2))\right|\\
 &\le C_7e^{-m_1(t-2)}\sum_{a,b}
       \|Q_a\|_{\rm GJS}\|Q_b\|_{\rm GJS}\\
 &=C_7e^{2m_1}e^{-m_1t}\|Q\|_{\rm GJS}^2.
 \end{split}
\]
This decomposition is required because GJS Theorem 1.1.7 is stated for localized
monomials; applying it directly to a general $Q$ would be invalid.

\LS{2}{5}{Let
$\psi_Q=Q\Om_g-\ip{\Om_g}{Q\Om_g}\Om_g$.  Then}
\[
 \ip{\psi_Q}{e^{-tH(g)}\psi_Q}
 \leq\widehat C_Qe^{-m_1t}\qquad(t\geq0),                          \tag{6.8}
\]
where
\[
 \widehat C_Q=\max\{C_7e^{2m_1}\|Q\|_{\rm GJS}^2,
                     e^{3m_1}\|\psi_Q\|^2\}.
\]

\Because For $t>2$, expand the connected expression, use (6.5), then (6.7),
and use (6.2) for both one-point terms.  Explicitly, if $a_T$ is the connected
Euclidean difference and $b_T$ the product of one-point functions, then
\[
 \liminf_T(a_T+b_T)\leq\sup_T|a_T|+\lim_Tb_T.
\]
For $0\leq t\leq3$, positivity of $H(g)$ gives
$\ip{\psi_Q}{e^{-tH(g)}\psi_Q}\leq\|\psi_Q\|^2
\leq e^{3m_1}\|\psi_Q\|^2e^{-m_1t}$.

\LS{2}{6}{This proves $\langle1\rangle6$.}\QedStep

\LS{1}{7}{Prove polynomial vectors are dense.}

Define $d\mu_g=|\Om_g|^2d\phi_0$.  Since $\Om_g>0$ almost everywhere,
multiplication by $\Om_g$ is a unitary map from $L^2(d\mu_g)$ to $\Hh$.
It therefore suffices to prove that cylinder polynomials are dense in $L^2(d\mu_g)$.

\LS{2}{1}{Every real cylinder variable $X_f=\varphi_0(f)$ has a nonzero
exponential-moment radius.}

\Because Let $f$ have support in an interval of length $L$, set
$K'=K_2^2N_I^{1/2}\|f\|_2$, and apply (6.1) to $Q=X_f^n$.  Then
\[
 \E_{\mu_g}[X_f^{2n}]
 \leq C_8K_1(K')^{2n}(2n)!.                                      \tag{7.1}
\]
Since $e^{s|x|}\leq2\cosh(sx)$,
\[
 \E_{\mu_g}e^{s|X_f|}
 \leq2\max\{1,C_8K_1\}\sum_{n=0}^\infty(sK')^{2n}<\infty
 \quad(0\leq s<K'^{-1}).                                         \tag{7.2}
\]
The factor $K_1$ in (7.2), omitted in the manuscript, affects only the prefactor,
not the radius.

\LS{2}{2}{If $\zeta\in L^2(d\mu_g)$ is orthogonal to all cylinder polynomials,
then its conditional expectation on every finite cylinder sigma-algebra is zero.}

\Because Fix real $e_1,\dots,e_k$ and define
\[
 F(z_1,\dots,z_k)=
 \E_{\mu_g}\left[\overline\zeta\exp\left(\sum_{j=1}^kz_jX_{e_j}\right)\right].
\]
By Cauchy--Schwarz, H\"older, and (7.2), $F$ is analytic in a nonempty
polystrip about $i\R^k$.  Differentiation under the integral is justified by the
same exponential dominator.  Every derivative at zero is an inner product with a
polynomial, hence is zero.  Repeated application of the one-variable identity theorem
gives $F=0$ throughout the connected polystrip, in particular on $i\R^k$.
Thus the finite complex measure
\[
 B\longmapsto\E_{\mu_g}\big[\overline\zeta\,
 \mathbf1_{\{(X_{e_1},\ldots,X_{e_k})\in B\}}\big]
\]
has zero Fourier transform.  Uniqueness of Fourier transforms of finite measures gives
$\E_{\mu_g}[\zeta\mid\sigma(X_{e_1},\dots,X_{e_k})]=0$.

\LS{2}{3}{A countable increasing family of these finite sigma-algebras exhausts the
cylinder sigma-algebra.}

\Because Choose $(e_j)$ dense in real $L^2(\R)$.  The free covariance
$C=(2\sqrt{-\partial_x^2+m_0^2})^{-1}$ is bounded, so
\[
 \E_{\phi_0}|X_f-X_e|^2=\ip{f-e}{C(f-e)}\leq\|C\|\|f-e\|_2^2.
\]
For an approximating sequence, a subsequence converges $\phi_0$-almost surely and
hence $\mu_g$-almost surely because $\mu_g\ll\phi_0$.  Therefore every $X_f$ is
measurable with respect to $\bigvee_k\sigma(X_{e_1},\dots,X_{e_k})$ modulo null sets.

\LS{2}{4}{The martingale convergence theorem now gives $\zeta=0$.}

\Because The conditional expectations in $\langle2\rangle2$ are all zero and converge
in $L^2$ to the conditional expectation on the exhausted cylinder sigma-algebra, which
is $\zeta$ itself.  Hence cylinder polynomials, and therefore
$\Span\{Q\Om_g\}$, are dense.

\LS{2}{5}{This proves $\langle1\rangle7$.}\QedStep

\LS{1}{8}{Deduce the physical Hamiltonian gap (G).}

Fix $0\leq\beta<m_1$ and let $P_\beta=\mathbf1_{[0,\beta]}(H(g))$.

\LS{2}{1}{For every centered polynomial vector $\psi_Q$, $P_\beta\psi_Q=0$.}

\Because The spectral theorem and (6.8) give, for every $t\geq0$,
\[
 e^{-\beta t}\|P_\beta\psi_Q\|^2
 \leq\ip{\psi_Q}{e^{-tH(g)}\psi_Q}
 \leq\widehat C_Qe^{-m_1t}.
\]
Multiply by $e^{\beta t}$ and let $t\to\infty$.  Since $m_1-\beta>0$,
$\|P_\beta\psi_Q\|=0$.

\LS{2}{2}{$P_\beta\Om_g=\Om_g$.}

\Because $H(g)\Om_g=0$ and $0\in[0,\beta]$.

\LS{2}{3}{$P_\beta=|\Om_g\rangle\langle\Om_g|$.}

\Because By $\langle1\rangle7$, the complex span of $\Om_g$ and the centered
polynomial vectors is dense.  Steps $\langle2\rangle1$--$\langle2\rangle2$ show that
$P_\beta$ maps this dense subspace into $\C\Om_g$.  Continuity of the bounded
projection extends this containment to all of $\Hh$, while
$P_\beta\Om_g=\Om_g$ shows its range is exactly that line.

\LS{2}{4}{There is no spectrum in $(0,m_1)$.}

\Because If an open interval $J\Subset(0,m_1)$ had a nonzero spectral projection,
choose $\beta<m_1$ above its right endpoint.  Then
$\mathbf1_J(H(g))\leq P_\beta-|\Om_g\rangle\langle\Om_g|=0$, a contradiction.
Thus $\spec H(g)\subseteq\{0\}\cup[m_1,\infty)$; simplicity was imported in
Assumption \ref{ass:cutoff}.

\LS{2}{5}{This proves (G), hence Theorem \ref{thm:EG}.}\QedStep

\section{Audit map back to the manuscript}

\begin{longtable}{P{1.8cm}P{3.0cm}P{4.5cm}P{5.0cm}}
\toprule
Location & Defect & Why the written step fails & Exact correction above\\
\midrule
\endhead
lines 305, 536 & wrong source & Simon III.12 is a projection-product estimate, not FKN. &
Assumption \ref{ass:fkn}; use Simon III.6 and V.14--V.15 or GRS II.14, II.16.\\
line 352 & positivity hypothesis unstated & The slab lemma requires each $F_i\geq0$, not merely $F_1F_2\geq0$. &
$\langle1\rangle4\langle2\rangle5$: $F_1=F_2=|V(h)|$.\\
lines 376--377 & false cell count & A length-$L$ interval can meet more than $L+1$ unit cells. &
$\langle1\rangle5$: use $N_I\leq\lfloor L\rfloor+2\leq L+2$.\\
lines 404--420 & incomplete unbounded passage & Signed observables are outside the stated nonnegative FKN lemma. &
$\langle1\rangle2\langle2\rangle2$--$\langle2\rangle3$ and
$\langle1\rangle6$: signed truncations plus an explicit $L^2$ dominator.\\
line 412 & false for $s<0$ & $\|\psi_{T+s}\|\leq\|\psi_T\|$ holds only for $s\geq0$. &
$\langle1\rangle6\langle2\rangle2$: split the two signs and use the bound $c_g^{-1}$.\\
line 455 & missing factor & $C_8K_1$ cannot disappear from the moment series. &
$\langle1\rangle7\langle2\rangle1$: keep it outside the geometric series.\\
line 484 & missing range restriction & $P_{[0,\beta]}\Om_g=\Om_g$ is false for $\beta<0$. &
$\langle1\rangle8$: state $0\leq\beta<m_1$.\\
standing assumption & scaling suppressed & The source smallness conditions are dimensional until scaling is written. &
$\langle1\rangle1$: exact dilation and definition of $\varepsilon_{\mathscr P}$.\\
\bottomrule
\end{longtable}

\section{Logical conclusion}

The argument establishes (M) and (G) for every cutoff $g$, with constants independent
of $g$, conditional on the enumerated construction facts and the GJS estimates.  None of
the identified defects forces a change in the claimed gap $m_1$ or in the qualitative
local-moment conclusion; each is repaired explicitly above.  The proof does not establish
the GJS estimates, cutoff-Hamiltonian self-adjointness, ground-state simplicity, or FKN;
those remain imported theorems and must not be hidden by the phrase ``conditional only on
GJS Theorems 1.1.7--1.1.8.''

\end{document}
